\clearpage\SourceLeaf{165}%
\pdfbookmark[1]{Summatio quarumdam serierum singularium. Review 1808 September 19}{notice-1808-09-19}
\begin{center}
\rule{.98\linewidth}{.3pt}\par
Göttingische Gelehrte Anzeigen.\quad 1808 September 19.\par
\rule{.98\linewidth}{.3pt}
\end{center}

A lecture submitted by Professor \textsc{Gauss} to the Royal Society of Sciences,
\begin{center}\emph{Summatio quarumdam serierum singularium},\end{center}
aims to develop more fully and in greater generality a remarkable investigation concerning the division of the circle, whose foundation had already been laid in the \emph{Disquisitionibus Arithmeticis}, to provide it with complete proofs, and to show its unexpected connection with other important truths. Let $n$ denote a prime, $k$ any integer not divisible by $n$, and $\omega$ the arc $\frac1n360^\circ$. Among the numbers $1,2,3,4,\dotsfour n-1$, let the distinct quadratic residues of $n$ be represented by $a,a',a''$, and so on, while the numbers remaining after these have been excluded, or the quadratic nonresidues of $n$, are represented by $b,b',b''$, and so on. It is proved in the cited work, Art.~356, that, when $n$ has the form $4n+1$,
\[
\left.\begin{aligned}
&\cos ak\omega+\cos a'k\omega+\cos a''k\omega+\text{etc.}\\
&-\cos bk\omega-\cos b'k\omega-\cos b''k\omega-\text{etc.}
\end{aligned}\right\}=\pm\surd n.
\]
and
\[
\left.\begin{aligned}
&\sin ak\omega+\sin a'k\omega+\sin a''k\omega+\text{etc.}\\
&-\sin bk\omega-\sin b'k\omega-\sin b''k\omega-\text{etc.}
\end{aligned}\right\}=0
\]
\SourceLeaf{166}%
whereas, when $n$ has the form $4m+3$, the sum of the first series becomes $=0$ and that of the second $=\pm\surd n$. The sign to be placed before the radical depends on the value of the number $k$, or rather on its relation to $n$, and is easily determined for \emph{all} values of $k$ at a given value of $n$ once it is determined for \emph{one}. For it can be shown that \emph{the same} sign applies to all values of $k$ that are quadratic residues of $n$, and the opposite sign to all those that are quadratic nonresidues of $n$. Since the investigation had already been carried this far in the cited work, and only the determination of the sign for any one value of $k$ remained, one would have thought that, with the main matter settled, this further determination could easily be supplied. This seemed all the more likely because examining examples immediately gives an extremely simple result: for $k=1$, or for all values that are quadratic residues of $n$, the radical in the formulas above must always be taken \emph{positively}. Yet in seeking a proof of this observation we encounter wholly unexpected difficulties, and the procedure that led so satisfactorily to the absolute value of those series proves entirely inadequate when the complete determination of the signs is required. The \emph{metaphysical} reason for this phenomenon (to use the expression customary among French mathematicians) must be sought in the fact that, in dividing the circle, analysis makes no distinction between the arcs $\omega,2\omega,3\omega\ldots(n-1)\omega$, but treats them all alike. As this gives the investigation a new interest, Professor \textsc{Gauss} found in it, as it were, an invitation to leave nothing untried in overcoming the difficulty. Only after many and varied unsuccessful attempts did he succeed, by a route remarkable in itself. He starts from the summation of certain series whose terms have the following form:
\[
\frac{(1-x^m)(1-x^{m-1})(1-x^{m-2})\ldots(1-x^{m-\mu+1})}{(1-x)\ (1-xx)\ (1-x^3)\dotsfour(1-x^\mu)}
\]
If, for brevity, we denote such a function by $(m,\mu)$, which, as the paper shows, is always a \emph{polynomial} in $x$, then the series
\SourceLeaf{167}%
\[
1-(m,1)+(m,2)-(m,3)+\text{etc.}
\]
\[
1+x^{\frac12}(m,1)+x(m,2)+x^{\frac32}(m,3)+\text{etc.}
\]
terminate after the term in position $m+1$, provided that $m$ denotes a positive integer. For even values of $m$ the sum of the first series becomes
\[
=(1-x)(1-x^3)(1-x^5)\ldots(1-x^{m-1})
\]
and $=0$ for odd values of $m$, whereas the sum of the second series always becomes
\[
=(1+x^{\frac12})(1+x)(1+x^{\frac32})\ldots(1+x^{\frac12m})
\]
For fractional and negative values of $m$, too, the summation of these series leads to interesting results, though they are not needed for the present purpose. We content ourselves with giving just one of them here. The infinite series
\[
1+x+x^3+x^6+x^{10}+\text{etc.}
\]
whose exponents are the triangular numbers is the product of the factors
\[
\frac{1-xx}{1-x}\times\frac{1-x^4}{1-x^3}\times\frac{1-x^6}{1-x^5}\times\frac{1-x^8}{1-x^7}\ \text{etc.}
\]
or, if preferred, of
\[
(1+x)^2(1+xx)^2(1+x^3)^2(1+x^4)^2\ \text{etc.}
\]
multiplied by
\[
(1-x)(1-xx)1-x^3)(1-x^4)\ \text{etc.}
\]

An account of how these summations are applied to the main subject would take us too far here. We may refer readers to the paper itself all the more readily since it will soon appear in print. The summations quoted above are only a special application of the summation of the following series:
\[
1+\cos k\omega+\cos4k\omega+\cos9k\omega+\text{etc.}+\cos(n-1)^2k\omega=T
\]
\[
\sin k\omega+\sin4k\omega+\sin9k\omega+\text{etc.}+\sin(n-1)^2k\omega=U
\]
which the paper gives for all values of $k$, and without the restriction
\SourceLeaf{168}%
that $n$ be prime. It is shown that
\[
T=\pm\surd n,\quad T=\pm\surd n,\quad T=0,\quad T=0
\]
and
\[
U=\pm\surd n,\quad U=0,\quad U=0,\quad U=\pm\surd n
\]
according as $n$ has the forms $4m,4m+1,4m+2,4m+3$, respectively. Here again the sign of the radical depends on $k$. Its determination requires many individual cases to be distinguished, and is developed and proved by two different routes in such a way that nothing will remain to be desired. Comparing these two routes also leads to the following very remarkable theorem: Let $n$ be the product of any number of distinct odd primes $a,b,c,d$, and so on, of which altogether $\mu$ have the form $4m+3$. Suppose further that among those factors there are altogether $\nu$ for each of which the product of the others (respectively $\frac na,\frac nb,\frac nc,\frac nd$, and so on) is a quadratic nonresidue. Then $\nu$ will be even whenever $\mu$ has the form $4m$ or $4m+1$, but odd whenever $\mu$ has the form $4m+2$ or $4m+3$. The known fundamental theorem concerning quadratic residues is only a special case of this theorem; conversely, the theorem just stated is easily derived from that fundamental theorem. Thus these investigations also provide a fourth proof of this important theorem, which the author first proved by two entirely different routes in the \emph{Disquisitionibus Arithmeticis}, and recently by a third, equally different route in a separate paper.
\BookRule{.28\linewidth}
